\section{The Gaussian central term for an exact physical singleton}
\label{sec:stationary-gaussian-complement}
We now evaluate the central part of
\eqref{eq:stationary-absolute-inversion}. This calculation takes place
in the actual three frequency variables $(u_1,u_2,b)$ of a prescribed
physical collision count. It is not a trace of a four-dimensional
integrated estimate. Together with the preceding uniform time tail it
reduces the original stationary singleton probability to one finite,
signed middle-frequency integral at a polynomial bandwidth.

Define the bounded centered collision record and its covariance by
\begin{equation}\label{eq:physical-three-collision-record}
 f^{\rm c}_R=(\kappa_{R,1},\kappa_{R,2},\tau_R-\bar\tau_R)
       =P_Rh_R,\qquad
 P_R(w_1,w_2,a,b)=(w_1,w_2,b-\bar\tau_Ra),\qquad
 \Sigma_R=P_R\Gamma_RP_R^{\mathsf T}.
\end{equation}
The map $P_R$ has full row rank with uniformly bounded coefficients.
Thus $\Sigma_R$ is uniformly positive definite by Theorem E. Its
Green--Kubo formula and its smoothed version follow by applying $P_R$
to the collision formulas; no new correlation theorem is needed.

\subsection{A three-dimensional endpoint-weighted central estimate}
\begin{lemma}[Two roof weights at the collision endpoints]
\label{lem:two-roof-central-comparison}
For $m\ge2$, uniformly in $R$,
\begin{align}
 &\int_{|v|\le2m^{1/200}}\left|
 \int_M\tau_R(x)\tau_R(T_R^mx)
       e^{iv\cdot\sum_{j<m}f^{\rm c}_R(T_R^jx)/\sqrt m}\,d\nu(x)
       -\bar\tau_R^2e^{-v^{\mathsf T}\Sigma_Rv/2}
 \right|\,dv\notag\\
 &\hspace{24mm}\le C m^{-11/700}\sqrt{\log(2+m)}.
 \label{eq:two-roof-central-comparison}
\end{align}
The integration is with respect to three-dimensional Lebesgue measure.
\end{lemma}
\begin{proof}
Put $h_\delta=\mathsf S_\delta h_R$ and
$\tau_\delta=\mathsf S_\delta\tau_R$, and let
$z=P_R^{\mathsf T}v/\sqrt m$. The collision transfer convention gives
exactly the smoothed pairing
\[
 \ell_j\bigl(M_{\tau_\delta}\mathcal L_{R,\delta,z}^{m}
                                        (\tau_\delta\nu)\bigr).
\]
The initial vector and terminal functional each have norm
$O(\delta^{-2})$. The spectral splitting of
Lemma~\ref{lem:smooth-collision-expansion} therefore has complementary
error $O(\delta^{-4}\rho^m)$. Its principal amplitude at $z=0$ is
$\bar\tau_R^2$, because smoothing preserves the mean. The difference
of that amplitude from its value at zero is bounded by
$C\delta^{-6}|v|/\sqrt m$: two endpoint norms and one projection
difference have all been paid.

Write $\ell_\delta=1+|\log\delta|$. The same spectral lemma gives
\[
 m\log\lambda_{R,\delta}(z)
 =-\tfrac12v^{\mathsf T}P_R\Gamma_{R,\delta}P_R^{\mathsf T}v
                           +O(\delta^{-6}|v|^3/\sqrt m).
\]
When the displayed cubic remainder and $|z|\delta^{-2}$ are small,
positive semidefiniteness of the smoothed covariance bounds the
principal power by a fixed constant. The covariance replacement costs
$C|v|^2\delta\ell_\delta$.

Replacing either endpoint $\tau_\delta$ by $\tau_R$ costs $C\delta$,
by invariance, the $L^1$ smoothing error and the uniformly bounded
other factors. The product of the two smoothed roofs is bounded
pointwise independently of $m,\delta$. The small-BV variance bound
\eqref{eq:small-bv-variance}, applied to each component of
$P_R(h_R-h_\delta)$, consequently bounds the unsmoothing of the phase
by $C|v|\sqrt{\delta\ell_\delta}$. The resulting pointwise error is
\[
 C\left[\delta+|v|\sqrt{\delta\ell_\delta}
 +|v|^2\delta\ell_\delta
 +m^{-1/2}\delta^{-6}(|v|+|v|^3)+\delta^{-4}\rho^m\right].
\]
No long pullback of an endpoint weight is estimated in BV.

Take $\delta=m^{-1/14}/4$ and integrate on the three-dimensional
ball of radius $2m^{1/200}$. In the order just displayed, the five
polynomial decay exponents are
\[
 \frac{79}{1400},\qquad \frac{11}{700},\qquad
 \frac{13}{280},\qquad \frac9{175},\qquad \frac{29}{700}.
\]
The second term is the largest; the extra logarithm in the covariance
term is absorbed by its strict power margin. The spectral smallness
requirements follow from
\[
 \frac12-\frac1{200}-\frac17>0,\qquad
 \frac12-\frac3{200}-\frac37>0.
\]
The complementary term is exponentially small. Increasing $C$ for
finitely many initial $m$ proves \eqref{eq:two-roof-central-comparison}.
\end{proof}

Set $r_m=2m^{-99/200}$, so that $\sqrt m\,r_m=2m^{1/200}$.
In the neighborhood $|(u,b)|\le r_m$, the bounded cell labels and ages
in Lemma~\ref{lem:finite-flight-cell-partition} give
\begin{equation}\label{eq:flight-amplitude-central-approximation}
 |A_R^\pm(x;u,b)-\tau_R(x)|\le C|(u,b)|.
\end{equation}
This estimate uses the exact cell labels. It does not require the
amplitudes themselves to have a uniform BV norm at every frequency.
For $m$ sufficiently large, this ball is inside the chosen torus chart.
Define its inverse contribution by
\begin{equation}\label{eq:stationary-central-inverse}
 G^{\rm cen}_{m,R}(k,t)=\frac1{(2\pi)^3\bar\tau_R}
 \int_{|(u,b)|\le r_m}
    e^{-i(u\cdot k+bt)}\mathcal F_{m,R}(u,b)\,du\,db.
\end{equation}

\begin{theorem}[Evaluated Gaussian central term]
\label{thm:stationary-singleton-central-term}
Writing
\[
 \zeta_{m,R}(k,t)=\left(\frac{k}{\sqrt m},
                             \frac{t-m\bar\tau_R}{\sqrt m}\right),
\]
one has, uniformly in $R,k,t$ and all sufficiently large $m$,
\begin{equation}\label{eq:stationary-singleton-central-term}
 \left|m^{3/2}G^{\rm cen}_{m,R}(k,t)
             -\bar\tau_R g_{\Sigma_R}(\zeta_{m,R}(k,t))\right|
                \le C m^{-11/700}\sqrt{\log(2+m)}.
\end{equation}
Here $g_{\Sigma}$ is the normalized Gaussian density on $\R^3$.
\end{theorem}
\begin{proof}
Remove $e^{ibm\bar\tau_R}$ from
\eqref{eq:stationary-physical-transform}. Substitution
$(u,b)=v/\sqrt m$ in \eqref{eq:stationary-central-inverse} gives
an inverse on the actual three-dimensional ball. By
\eqref{eq:flight-amplitude-central-approximation}, replacement of
the two amplitudes by the endpoint roofs costs at most
\[
 C m^{-1/2}\int_{|v|\le2m^{1/200}}|v|\,dv
                         \le C m^{-12/25}
\]
after multiplication by $m^{3/2}$. Apply
Lemma~\ref{lem:two-roof-central-comparison} to the resulting transform.
The error in an inverse is bounded by its integrated transform error,
uniformly in the target. Uniform ellipticity gives a Gaussian tail
$Ce^{-c m^{1/100}}$ outside this ball. The full Gaussian inverse is
$g_{\Sigma_R}(\zeta)$ with factor $\bar\tau_R^2/\bar\tau_R$.
This proves the normalization and the asserted bound.
\end{proof}

\subsection{One finite signed complement for the original physical event}
For $B\ge1$ and sufficiently large $m$ put
\begin{equation}\label{eq:stationary-middle-complement}
 \mathfrak M_{m,R,B}(k,t)
 =\frac1{(2\pi)^3\bar\tau_R}
 \int_{\substack{u\in[-\pi,\pi]^2,\ |b|\le B\\ |(u,b)|>r_m}}
 e^{-i(u\cdot k+bt)}\mathcal F_{m,R}(u,b)\,db\,du.
\end{equation}
This is a signed inverse evaluated at the target, not an integral of
an absolute value. Its domain contains all noncentral spatial torus
frequencies and all remaining time frequencies below $B$.

\begin{theorem}[A pointwise Gaussian-plus-complement identity]
\label{thm:physical-singleton-reduction}
For every $B\ge1$ and all sufficiently large $m$,
\begin{align}
 &\sup_{R,k,t}\big|m^{3/2}p_{m,R}(k,t)
  -\bar\tau_Rg_{\Sigma_R}(\zeta_{m,R}(k,t))
  -m^{3/2}\mathfrak M_{m,R,B}(k,t)\big|\notag\\
 &\hspace{18mm}\le C\left(m^{-11/700}\sqrt{\log(2+m)}+\frac{m^{3/2}}B\right).
 \label{eq:physical-singleton-local-ledger}
\end{align}
There is no discarded-density term in this identity. For
$B_m=m^{P+3/2}$, its last term is $O(m^{-P})$.

More intrinsically, put
\[
 \xi_{t,m,R}(k)=\left(\frac{k}{\sqrt t},
                              \frac{m-t/\bar\tau_R}{\sqrt t}\right).
\]
For every fixed $M<\infty$, uniformly in $R$ and $k,m,t$ satisfying
$t\to\infty$ and $|\xi_{t,m,R}(k)|\le M$,
\begin{align}
 t^{3/2}p_{m,R}(k,t)
 &=g_{\mathcal V_R}(\xi_{t,m,R}(k))
       +t^{3/2}\mathfrak M_{m,R,B_m}(k,t)\notag\\
 &\quad+O_M\left(t^{-11/700}\sqrt{\log(2+t)}+t^{-P}\right).
 \label{eq:physical-singleton-intrinsic-ledger}
\end{align}
Thus the Gaussian denominator asymptotic for these original physical
singletons is equivalent to the vanishing of the displayed normalized
signed complement on each such central set.
\end{theorem}
\begin{proof}
Split the absolutely convergent integral in
\eqref{eq:stationary-absolute-inversion} into its central ball, the
finite middle region \eqref{eq:stationary-middle-complement}, and
$|b|>B$. Theorem~\ref{thm:stationary-singleton-central-term} evaluates
the first part and Theorem~\ref{thm:stationary-polynomial-band} bounds
the third. This gives \eqref{eq:physical-singleton-local-ledger} for
the actual event, pointwise in observation time. No essential supremum
of a discarded return density has entered the argument.

Let $D_R^{\rm clk}=\operatorname{diag}(1,1,-1/\bar\tau_R)$. The
physical covariance defined earlier satisfies exactly
\begin{equation}\label{eq:physical-covariance-jacobian}
 \mathcal V_R=\bar\tau_R^{-1}
              D_R^{\rm clk}\Sigma_R(D_R^{\rm clk})^{\mathsf T},\qquad
 \det\mathcal V_R=\det\Sigma_R/\bar\tau_R^5.
\end{equation}
Indeed $D_R^{\rm clk}P_R$ is the map $B_R$ of the physical-time
functional theorem. If $|\xi|\le M$, then
$m=t/\bar\tau_R+O_M(\sqrt t)$ uniformly. Set
$A_R=\sqrt{\bar\tau_R}\operatorname{diag}(1,1,-\bar\tau_R)$.
Its absolute determinant is $\bar\tau_R^{5/2}$ and
$A_R\mathcal V_RA_R^{\mathsf T}=\Sigma_R$. Also
\[
 \zeta_{m,R}=A_R\xi+O_M(t^{-1/2}),\qquad
 \bar\tau_R(t/m)^{3/2}=\bar\tau_R^{5/2}+O_M(t^{-1/2}).
\]
Uniform Gaussian derivative bounds on compact sets therefore show
\[
 \bar\tau_R(t/m)^{3/2}g_{\Sigma_R}(\zeta_{m,R})
                  =g_{\mathcal V_R}(\xi)+O_M(t^{-1/2}).
\]
Multiplying \eqref{eq:physical-singleton-local-ledger} by $(t/m)^{3/2}$
proves \eqref{eq:physical-singleton-intrinsic-ledger}. Since the
remaining explicit error tends to zero, subtraction proves both
directions of the asserted equivalence. A signed complement may
vanish even if its absolute transform integral is not small.
\end{proof}

\begin{corollary}[What a complement estimate would imply]
\label{cor:stationary-singleton-denominator-criterion}
If the normalized middle term in
\eqref{eq:physical-singleton-intrinsic-ledger} tends to zero uniformly
on a fixed central set, then there
\[
 \mathbb P_R\{W_R(t)=k,C_R(t)=m\}
           =t^{-3/2}g_{\mathcal V_R}(\xi_{t,m,R}(k))(1+o(1)),
\]
and this denominator is uniformly bounded below by $c_Mt^{-3/2}$.
For any $P>0$, Theorem~\ref{thm:physical-time-selector-reconstruction}
with $B_m=m^{P+3/2}$ then approximates its conditional initial law by
the positive finite-band likelihood with total variation error
$O_M(t^{-P})$. A selected event needs its own selected denominator;
the unselected Gaussian formula does not establish it.
\end{corollary}
\begin{proof}
Uniform ellipticity makes the Gaussian density uniformly positive on
the compact central set. Use
Theorem~\ref{thm:physical-singleton-reduction} and divide the source
error in Corollary~\ref{cor:direct-physical-posterior} by this lower
bound. The final sentence specifies the normalization required after
multiplication by a selector.
\end{proof}

\subsection{Relation to the return-density and suspension theories}
The cell-displacement local theorem of \cite{SV} and the suspension
mixing-local-limit framework of \cite[Definition 2.2 and Section 2.4]{DN}
are natural comparisons. In the latter framework a local theorem on
the base is an input to the suspension argument. We have not imported
that input for the moving-section four-coordinate record. The present
calculation instead keeps one prescribed physical collision count and
the complete cell corrections, integrates the true flight ages, and
obtains a uniform pointwise far-time-frequency estimate directly.
Its two single-flight amplitudes and their inverse-square product
are the mechanism behind the polynomial bandwidth. They do not by
themselves evaluate the middle-frequency collision correlation.

For reference the three norms and the completed steps can be recorded
without identifying their endpoints:
\begin{center}
\small
\renewcommand{\arraystretch}{1.15}
\begin{tabular}{@{}p{0.23\textwidth}p{0.37\textwidth}p{0.32\textwidth}@{}}
\hline
Object & Proved estimate & Remaining local task\\
\hline
Preceding-section return law & Theorem W: global mixed $L^1$ and
source-TV approximation; $B_n$ may be $\exp(O(n\log n))$ & The common
pointwise return correction and its finite complement.\\
Actual physical singleton & Theorem X: a pointwise $O(B^{-1})$ tail,
a Gaussian central term and a polynomial-band exact ledger & The
normalized signed middle term in
\eqref{eq:physical-singleton-intrinsic-ledger}.\\
Physical weighted event & Absolute source-TV error $E_B^{\rm phys}$
for every bounded selector & Its evaluated weighted formula and positive
selected denominator before relative conditioning.\\
\hline
\end{tabular}
\end{center}
The local physical remainder is now one explicitly specified integral,
with no discarded-density or event-replacement error. This is a direct
analysis of the original physical event, not a proof that it coincides
with a preceding-section event. The full four-coordinate return LLT
and the microscopic Gaussian denominator remain the stated targets;
neither is relabeled as a consequence of the new tail estimate.
